\section{Related work}
\label{sec:related}

\paragraph{Batch size and gradient noise.} The gradient noise scale predicts the batch size beyond
which data parallelism stops buying speed, and does so from statistics measurable during training
\citep{mccandlish2018empirical}; the empirical picture across tasks and optimisers is mapped in
\citet{shallue2019measuring}. Both concern a single level of sampling. RLVR samples prompts and then
samples responses within each prompt, and the resulting two-level noise scale, the split it implies,
and the fact that the response-level term is readable off the reward histogram, have no counterpart
in that literature. The variance decomposition itself is the classical two-stage sampling result
\citep{cochran1977sampling}, applied here to policy gradients.

\paragraph{Hyperparameter transfer.} $\mu$P makes learning rates transfer across width by fixing what
a step does to the network's functions rather than to its parameters
\citep{yang2021tensor, yang2022mutransfer}. This paper takes the same stance one level up: fix what
a step does to the policy's \emph{behaviour}, measured as drift in nats per token, and the recipe
transfers across batch size. We find no comparable prescription for RL post-training in the
literature; the closest work reports empirical scaling behaviour and fitted power laws
\citep{tan2025scalingbehaviors, khatri2025artofscaling, ghosh2026predictable} rather than a
quantity to hold fixed.

\paragraph{Estimators for RLVR.} RLOO \citep{ahmadian2024rloo, williams1992reinforce}, GRPO
\citep{shao2024deepseekmath}, its unnormalised correction \citep{liu2025drgrpo}, DAPO
\citep{yu2025dapo} and sequence-level variants \citep{zheng2025gspo} are usually compared by
benchmark deltas. Proposition~\ref{prop:mean} shows that, for binary rewards, the members of this
family differ only through a difficulty weight $\lambda(p, G)$, that a $p$-independent $\lambda$ is
absorbed exactly by the step size, and that standardisation is a change of objective rather than a
variance reduction -- a sharpening of the argument in \citet{liu2025drgrpo}. Recent analyses of why
normalisation helps \citep{ge2026normalization} and of what outcome-reward objectives cannot
simultaneously be \citep{ding2026impossibility} are complementary: they characterise the estimator,
while the question here is how many samples to give it and how far to step.

\paragraph{Allocating rollouts.} Several recent methods vary the number of responses per prompt
within a batch using heuristics tied to observed pass rates or gradient variance
\citep{nguyen2026allocation}, or adapt the batch size from a divergence signal
\citep{park2026batchscaling}, and one line of work argues for scaling the response count
specifically \citep{hu2025brorl}. These are allocation policies; what has been missing is the
quantity being optimised. Equation~\eqref{eq:gstar} supplies it in closed form and shows that the
answer is a ratio of two measurable traces, which also explains why heuristics keyed to pass rate
alone succeed on some corpora and not others: pass rate determines $\tau_w$ but says nothing about
$\tau_b$.

\paragraph{Trust regions and KL control.} Constraining how far a policy moves is old
\citep{schulman2015trpo, schulman2017ppo} and its modern LLM forms clip importance ratios or
penalise divergence from a reference. The drift used here is neither: it is the KL between
\emph{consecutive} policies, used as a measurement and a control target rather than as a penalty,
and its role in this paper is to supply the unit in which a recipe transfers.

\paragraph{Asynchrony and off-policy correction.} A parallel literature studies what happens when
the sampling policy and the training policy differ, whether through pipeline staleness
\citep{song2026staleness} or numerical mismatch between rollout and training engines. That work
asks how much mismatch a run tolerates. This paper asks a question that arises before any mismatch:
given that the sampler and the trainer agree, how should the sampling budget be divided, and how
large a step should the result support.
